\chapter{Prerequisites}
\label{ch:prereq}

\tbd

% Phase 0. The background a reader needs before Spark internals make sense:
% the relational model, the JVM, distributed-systems vocabulary, and how
% analytic data is stored on disk.

\section{The Relational Model and SQL}
\label{sec:prereq-sql}

\subsection{Relational algebra and SQL semantics}
\tbd

\subsection{Logical versus physical query plans}
\tbd

\subsection{Joins, aggregation, and window functions}
\tbd

\section{The JVM Execution Environment}
\label{sec:prereq-jvm}

Spark runs on the Java Virtual Machine: the driver and every executor is a JVM
process, and the engine's memory behaviour, its failure modes, and much of its
internal design follow directly from how the JVM manages memory and runs code.
This section covers the parts of that environment a reader needs before the
later chapters. The Spark-specific memory model built on top of it is
\Cref{sec:fund-memory}; the tuning knobs are \Cref{ch:tuning}.

\subsection{Bytecode and just-in-time compilation}
\tbd

% No source material yet in spark_memory_tuning.md; to be written alongside
% whole-stage code generation (Sec 3.2.5 / Sec 9.2).

\subsection{Heap layout and object representation}
\label{sec:prereq-heap}

The JVM manages a single contiguous region of memory, the \emph{heap}, from which
every object a program allocates is drawn. A Spark executor is one JVM process,
and the value of \texttt{spark.executor.memory} is precisely the maximum size of
that heap. Everything a task does --- decoding a Parquet page into rows, building
a hash table for a join, accumulating a group-by result --- competes for space in
it, and when it fills, the task must either spill to disk or fail.

Two properties of the heap drive most of Spark's decisions about memory. The
first is \emph{per-object overhead}. Every object carries a header --- a mark word
and a pointer to its class --- before any field is stored, and the allocator pads
every object to an eight-byte boundary.

\begin{lstlisting}[numbers=none,caption={Cost of a boxed integer on a 64-bit JVM},captionpos=b]
java.lang.Integer
  object header   12 bytes   (mark word + class pointer)
  int value        4 bytes
  ----------------------
  total           16 bytes   to carry 4 bytes of data
\end{lstlisting}

\noindent A row represented as an array of boxed Java objects (an
\texttt{Integer}, a \texttt{Double}, a \texttt{String} with its own header and
backing character array) can be several times larger in the heap than the same
row encoded as raw bytes. This is the amplification a reader sees when a
compressed Parquet file expands to many times its on-disk size once decoded into
Java objects, and it is why Spark's execution engine stores rows off the object
heap in a packed binary format rather than as JVM objects
(\Cref{sec:int-tungsten}).

The second property is that the heap is \emph{managed}: the program never frees
memory explicitly, and reclamation is the job of the garbage collector
(\Cref{sec:prereq-gc}). Spark subdivides the heap it is given into a small
reserved region, a fraction for the objects user code creates, and a fraction it
manages explicitly for execution and caching --- the unified memory model of
\Cref{sec:fund-memory}. That boundary is a Spark convention layered on one flat
JVM heap, not a JVM feature.

\subsection{Garbage collection models: Parallel, G1, ZGC}
\label{sec:prereq-gc}

Because the heap is managed, the JVM must periodically find objects that are no
longer reachable and reclaim their space. This work is the \emph{garbage
collector} (GC), and for a data-processing workload it is never free: the
collector competes with the application for CPU, and some of its phases are
\emph{stop-the-world}, meaning every application thread in the JVM is frozen until
the phase completes. On an executor a long stop-the-world pause stalls every
running task at once, and if it exceeds the heartbeat timeout the driver declares
the executor lost and re-runs its work elsewhere (\Cref{sec:dbg-failures}).

Collectors differ mainly in how they trade pause length against throughput and
CPU cost.

\begin{itemize}
  \item \textbf{Parallel GC} uses multiple threads but collects the whole heap in
        one stop-the-world pause. It gives the highest raw throughput when pauses
        do not matter, but pause time grows with heap size, so a large executor
        can freeze for seconds at a time.
  \item \textbf{G1} (Garbage-First) divides the heap into equal-sized regions and
        collects a subset of them per cycle, doing most of its marking
        concurrently with the application. Pause times are shorter and roughly
        bounded regardless of heap size, at the cost of more bookkeeping and CPU.
        It is the default at the heap sizes a Spark executor typically runs and
        the collector most Spark tuning assumes.
  \item \textbf{ZGC} does essentially all of its work concurrently and keeps
        pauses under a millisecond even on very large heaps, at a higher constant
        CPU and memory overhead. It is worth choosing only for executors with
        tens of gigabytes of heap where G1 pauses are still a problem.
\end{itemize}

The generational hypothesis --- that most objects die young --- underlies all
three: they collect a small ``young'' area cheaply and often, and the full-heap
collection rarely. A workload that holds a large amount of long-lived data in
cache violates that assumption and pushes the collector toward the expensive
path, which is one more reason to cache deliberately (\Cref{sec:qo-caching}).
Monitoring GC time is covered in \Cref{sec:dbg-metrics}; the executor JVM options
that select and tune a collector are in \Cref{ch:tuning}.

\subsection{Off-heap memory and the cost of serialization}
\label{sec:prereq-serialization}

Not all of a JVM process's memory is on the heap. The class metadata the JVM
loads, each thread's call stack, buffers for network and file I/O, memory
allocated by native libraries such as compression codecs, and any region the
application explicitly requests \emph{off-heap} all live outside the heap and
outside the garbage collector's control. On Kubernetes this is why an executor
pod is sized as \texttt{spark.executor.memory} plus an overhead: the heap is only
part of what the process touches, and if the total exceeds the pod's limit the
kernel kills it with no Java-level error (\Cref{sec:dbg-failures}).

Off-heap memory matters to Spark for a second reason: it is where the execution
engine puts data it wants to keep away from the garbage collector. A hash table
for a join, or a block of cached rows, held as one large off-heap byte buffer
instead of millions of small Java objects, is invisible to GC and costs nothing
to collect (\Cref{sec:int-tungsten}).

\emph{Serialization} is the conversion of an in-memory object graph into a flat
sequence of bytes, and \emph{deserialization} is its inverse. Spark serializes
whenever data leaves a JVM: across the network during a shuffle, to disk when a
partition spills or is cached in serialized form, and when a broadcast variable
is shipped to executors. The cost is both CPU and size, and the two Java options
differ sharply. The built-in Java serializer writes the full class name and every
field name and type alongside the values, repeated for every object, so most of
the output is self-description rather than data. \emph{Kryo} writes a compact
encoding --- a short numeric class identifier and the field values --- and
produces output several times smaller, which directly reduces shuffle volume,
spill size, and broadcast footprint (\Cref{sec:tune-serialization}).

DataFrame and SQL workloads sidestep both serializers for most operations: rows
already live in the engine's binary format (\Cref{sec:int-tungsten}), so
shuffling them is a memory copy with no object graph to walk. The Java-versus-Kryo
choice still governs RDD code, user-defined functions, and anything that puts
arbitrary Java objects into a shuffle or a broadcast.

\section{Distributed Systems Concepts}
\label{sec:prereq-distsys}

\subsection{Partitioning and data parallelism}
\tbd

\subsection{The shuffle as distributed data exchange}
\tbd

\subsection{Fault tolerance, lineage, and idempotency}
\tbd

\subsection{Consistency models and the CAP trade-off}
\tbd

\section{File Formats and Storage}
\label{sec:prereq-formats}

\subsection{Row-oriented versus columnar layouts}
\tbd

\subsection{Parquet: encoding, compression, and footer statistics}
\tbd

\subsection{ORC and Avro}
\tbd

\subsection{Object storage semantics (GCS and S3)}
\tbd

\section{Summary}
\tbd
